\documentclass[11pt,a4paper]{article}
\usepackage[T1]{fontenc}
\usepackage{lmodern,amsmath,amssymb,amsthm,booktabs,geometry,hyperref,xurl}
\geometry{margin=27mm}
\newtheorem{theorem}{Theorem}[section]
\newtheorem{lemma}[theorem]{Lemma}
\newtheorem{proposition}[theorem]{Proposition}
\newtheorem{definition}[theorem]{Definition}
\newcommand{\R}{\mathbb R}
\newcommand{\Q}{\mathbb Q}
\newcommand{\N}{\mathbb N}
\newcommand{\conv}{\operatorname{conv}}
\newcommand{\dd}{\sqrt{462}}
\newcommand{\file}[1]{{\normalfont\small\nolinkurl{#1}}}
\title{An explicit interval in the Markov spectrum below Hall's ray\\
\large A computer-assisted proof with a formally verified certificate checker}
\author{Koki Yamada}
\date{October 9, 2026}
\begin{document}
\maketitle
\begin{abstract}
We give a computer-assisted proof that the exact rational interval
$[4.52578,4.52754]$ is contained in the Markov spectrum strictly below the classical endpoint of
Hall's ray,
$c_F=(2221564096+283748\sqrt{462})/491993569$. A closed successor cover for continued-fraction tails avoiding
$31313$ is attached to the outward prefixes $(322,431)$, with central
digit $4$. Its finite certificate has 3,464,816 adopted rows. We explain
uniform verification over parameter boxes, passage from the finite cover
to actual infinite tails, and a bound at every noncentral position.
Lean verifies the implication from checker acceptance to the spectral
inclusion; compiled execution, rather than kernel reduction of the large
table, establishes concrete acceptance. Independent implementations and
reproduction instructions document the remaining execution trust boundary.
No assertion about the Lagrange spectrum or the unfinished $131$-forbidden
construction is made.
\end{abstract}

\section{Statement and scope}
For a sequence $\theta\in\N_{>0}^{\mathbb Z}$, put
\[
 \lambda_n(\theta)=\theta_n+[0;\theta_{n-1},\theta_{n-2},\ldots]
                         +[0;\theta_{n+1},\theta_{n+2},\ldots],\qquad
 m(\theta)=\sup_{n\in\mathbb Z}\lambda_n(\theta).
\]
We use the continued-fraction definition
$M=\{m(\theta):\theta\in\N_{>0}^{\mathbb Z},\ m(\theta)<\infty\}$
of the Markov spectrum. Its equivalence with the classical quadratic-form
definition is background theory; see Cusick--Flahive \cite{CF}.
Freiman's Hall-ray theorem \cite{Freiman} identifies
\[
 c_F=\frac{2221564096+283748\dd}{491993569}
\]
as the left endpoint of the maximal terminal ray. That classical
identification is used to interpret ``below Hall's ray''; it is not
reproved or included in the Lean development considered here.

\begin{theorem}\label{thm:main}
The computer-assisted verification described below establishes
\[
 \left[\frac{226289}{50000},\frac{226377}{50000}\right]
 =[4.52578,4.52754]\subset M\cap(-\infty,c_F).
\]
Consequently $(4.52578,4.52754)\subset\operatorname{int}_{\R}
(M\setminus[c_F,\infty))$. In particular $226333/50000=4.52666$ is an
interior point outside Hall's ray. The interval has width $11/6250$.
\end{theorem}
All terminating decimals in exact statements denote rational numbers.
The proof is supported by finite exact computations, not by sampling
spectral values. We state the finite obligations explicitly in
Sections~\ref{sec:graph} and \ref{sec:spectral}; the large obligation is
accepted by a compiled checker whose soundness is kernel checked.
The distinction is part of the statement of the verification method,
not an assertion of an unconditional closed Lean theorem.
The original release \cite{release} fixes the computational evidence.
The source snapshot inspected for this exposition is commit
\file{4b38497822c2088f9103a3c31aa4a1ac9a2a1462}.

\section{Continued fractions and the forbidden-word language}
For a finite word $U=(u_1,\ldots,u_h)$ of positive integers let
$F_U(x)=[0;u_1,\ldots,u_h+x]$, with $F_{\varnothing}(x)=x$.
Write $p_j/q_j=[0;u_1,\ldots,u_j]$, where
$p_{-1}=1,p_0=0,q_{-1}=0,q_0=1$ and
$p_j=u_jp_{j-1}+p_{j-2}$, $q_j=u_jq_{j-1}+q_{j-2}$.
Then
\begin{equation}\label{eq:mobius}
 F_U(x)=\frac{p_h+p_{h-1}x}{q_h+q_{h-1}x},\quad
 F_U(x)-F_U(y)=\frac{(-1)^h(x-y)}
 {(q_h+q_{h-1}x)(q_h+q_{h-1}y)}.
\end{equation}
These follow by multiplying the matrices of $x\mapsto1/(u+x)$;
the determinant identity is $p_{h-1}q_h-p_hq_{h-1}=(-1)^h$.
For nonempty $U$ set $r_U=q_{h-1}/q_h\in[0,1]$.
Equation~\eqref{eq:mobius} implies that the image of $[0,1]$ has width
at most $q_h^{-2}$. Since $q_h\ge F_{h+1}$, with
$F_1=F_2=1$, these widths tend to zero uniformly in words of length $h$.
This also proves convergence of the infinite continued fractions used here.

For a fixed word $U$ avoiding $31313$, define
\[
 K(U)=\{[0;U,w]: w\in\{1,2,3\}^{\N},\ Uw\text{ avoids }31313\}.
\]
The fixed word may contain $4$; only appended digits must be at most $3$.
Thus $K(431)$ is nonempty. Infinite extensions form a closed subset of the
compact product $\{1,2,3\}^{\N}$, and the uniform width estimate proves
continuity of their continued-fraction values. Hence $K(U)$ is compact.
Appending only $2$ always gives a legal extension.

The automaton state is the longest suffix which is a proper prefix of
$31313$. Label its five states $0,1,2,3,4$ by
$\varnothing,3,31,313,3131$, respectively. Its transition table is
\[
\begin{array}{c|ccc}
 s&1&2&3\\\hline
 0&0&0&1\\1&2&0&1\\2&0&0&3\\3&4&0&1\\4&0&0&\text{reject}
\end{array}
\]
A fixed digit $4$ resets the state to $0$. Write $T_s$ for the possible
values of legal infinite tails starting in state $s$.

\begin{lemma}\label{lem:extrema}
Writing $d=\sqrt{462}$, the endpoints $l_s=\min T_s$ and
$u_s=\max T_s$ are
\[
\begin{array}{c|cc}
 s&l_s&u_s\\\hline
0&(-29+2d)/53&-1+d/12\\
1&(-29+2d)/53&(-28+2d)/19\\
2&-1/2+d/28&-1+d/12\\
3&(-29+2d)/53&(-29+2d)/19\\
4&(-24+2d)/53&-1+d/12
\end{array}
\]
\end{lemma}
\begin{proof}
For every legal transition $s\xrightarrow{a}t$, direct substitution gives
\[
 l_s\le\frac1{a+u_t},\qquad \frac1{a+l_t}\le u_s.
\]
Denominators are positive; the inequalities can be checked by
cross multiplication and $d^2=462$, $21<d<22$.
These are exactly fourteen pairs of inequalities, one for every legal
state/digit pair. Applying them backwards along a finite legal prefix,
with any terminal point in $[l_t,u_t]$, gives a value in $[l_s,u_s]$.
Replacing the terminal point by any point of $[0,1]$ changes the value
by at most $q_h^{-2}$. Letting $h\to\infty$ proves $T_s\subset[l_s,u_s]$.
For attainment, at every state choose a legal digit realizing each of
\[
 l_s=\min_{s\xrightarrow{a}t}\frac1{a+u_t},\qquad
 u_s=\max_{s\xrightarrow{a}t}\frac1{a+l_t}.
\]
These equalities follow from the same finite substitutions.
Alternate lower and upper choices as the reciprocal reverses order.
The resulting infinite legal word has value the selected endpoint:
after $h$ digits its endpoint value differs from that word's value by
at most $q_h^{-2}$. This proves both attainment and the formulas.
\end{proof}
Let $L_U=\min K(U)$ and $H_U=\max K(U)$.
By parity in \eqref{eq:mobius}, they are the minimum and maximum of
$F_U(l_{\sigma(U)})$ and $F_U(u_{\sigma(U)})$.
For $0\le p\le q\le1$ define
\begin{equation}\label{eq:band}
 J_{p,q}(U,V)=[L_U+L_V+pW,\ L_U+L_V+qW],\qquad
 W=H_U+H_V-L_U-L_V.
\end{equation}
These are subintervals of a convex hull; membership in the actual Cantor
sum requires a separate argument.

\section{Closed covers and infinite realization}\label{sec:infinite}
Let $\xi_U$ be the endpoint $l_{\sigma(U)}$ or $u_{\sigma(U)}$ for which
$F_U(\xi_U)=L_U$. Define
\[
 D_U=\frac1{q_h^2(1+r_U\xi_U)^2},\qquad S(U,V)=D_V/D_U.
\]
In particular $1/4\le D_Uq_h^2\le1$.

\begin{lemma}[Closed successor-cover principle]\label{lem:cover}
Suppose a family of actual legal word pairs $(U,V)$ with assigned bands
$J_{p,q}(U,V)$ satisfies the following conditions. Each band is covered
by finitely many bands at extended word pairs $(Uu,Vv)$ in the same
family, where the added digits lie in $\{1,2,3\}$ and
$|u|+|v|>0$. Uniformly throughout the family,
$0<c\le D_V/D_U\le C<\infty$.
Then every assigned band is contained in $K(U)+K(V)$.
\end{lemma}
\begin{proof}
Fix a point $z$ of an assigned band. Choose successively a child band
containing $z$. This gives nested prefixes $U_n,V_n$, with total length
increasing by at least one at each step. At least one length tends to
infinity and its denominator tends to infinity. The scale of that side
tends to zero. If the other length were bounded, its nested prefixes
would eventually be constant, so its scale would be bounded below by
$1/(4q^2)>0$, contradicting the ratio bounds. Hence both lengths tend
to infinity. (Even if its anchors changed, the same lower bound applies.)
The compact sets $K(U_n)$ and $K(V_n)$ are nested and nonempty;
their diameters tend to zero by \eqref{eq:mobius}. Their intersections
are singletons $x\in K(U)$ and $y\in K(V)$.
Each band containing $z$ lies in
$[L_{U_n}+L_{V_n},H_{U_n}+H_{V_n}]$, whose diameter tends to zero.
Thus $z=x+y$. Every finite forbidden-word constraint persists in the
limit since an occurrence would already appear in a finite prefix.
This proves the claim. Notice that the auxiliary bands themselves need
not be nested.
\end{proof}

\section{The finite successor certificate}\label{sec:graph}
We reuse the saved \file{graph_wide} adoption table, rather than rerun
an exploratory search to choose a new family. A one-sided finite state
records the automaton state, prefix parity, and a shape box for $r_U$.
There are 22 such states and 484 ordered state pairs.
Shape boxes are obtained from reversed two-digit suffix maps applied to
$[1/4,4/5]$. This history interval is invariant under $x\mapsto1/(a+x)$
for $a=1,2,3$: the union of the three images lies in $[1/4,4/5]$.
The root shapes are checked separately, including the fixed $4$.
Ratios are partitioned by powers of $\kappa=51/50$.
A row consists of a state pair, ratio box, and band type $(p,q)$.
The fixed bitmap contains 3,464,816 adopted rows in 150,040 nonempty
cells. A cell groups rows with the same state pair and ratio box.

\subsection{Uniform endpoint comparisons}
The comparisons must hold for all real shapes and ratios in a box.
If $\xi$ is the anchor and $x,y$ are endpoint tails, then
\begin{equation}\label{eq:normalized}
 \frac{F_U(x)-F_U(y)}{D_U}
 =(-1)^{|U|}(x-y)
    \frac{1+r_U\xi}{1+r_Ux}\frac{1+r_U\xi}{1+r_Uy}.
\end{equation}
For fixed $\xi,x$, the ratio
$R(r)=(1+r\xi)/(1+rx)$ is monotone, since
$R'(r)=(\xi-x)/(1+rx)^2$.
Consequently its values at the shape-box endpoints bound it throughout
that box. Rational outward enclosures of the endpoint tails give bounds
on \eqref{eq:normalized}. Left differences are normalized by $D_U$;
right differences contribute with the factor $S=D_V/D_U$.

Endpoint comparisons preserve cancellations between repeated endpoint
labels. For example, the difference between two interpolated endpoints
has coefficients $1-p,p,-(1-q),-q$, whose sum is zero.
The checker verifies an exact identity
\[
 \sum_i c_i z_i=\sum_e w_e(z_{a_e}-z_{b_e}),\qquad w_e\ge0,
\]
by comparing the coefficient of every endpoint label. Each difference
on the right has a prepared rational enclosure. Multiplication by
$w_e$ and addition therefore enclose the left side. Two proposed
transports may be intersected to sharpen the enclosure; both identities
are checked. A nonnegative lower bound proves the requested endpoint
inequality. Proposing a transport is not a trusted step.

For extension $u$, the derivative multiplier is
\[
 \frac{D_{Uu}}{D_U}
 =\frac{(q_U+q_{U^-}\xi_U)^2}
        {(q_{Uu}+q_{(Uu)^-}\xi_{Uu})^2}.
\]
Its enclosure is validated by exact endpoint semantics and shape
transport. Hence a child ratio belongs to an enclosure of
$S(D_{Vv}/D_V)/(D_{Uu}/D_U)$. All selected band types must be adopted
in every destination bin of a certified contiguous range enclosing
that image. Exact equal nonzero multiplier tags allow the child ratio
to remain in the parent bin. Tags are checked as quadratic-field
identities, not inferred from overlapping numerical enclosures.
At a bin boundary it suffices to choose one verified bin containing
the actual ratio; no assumption that every touching bin is adopted
is needed for that choice.

\subsection{What is checked for each row}
A child vertex records legal extensions, a destination-bin range,
and selected band types. A merged child band is admitted only if its
whole relative-coordinate interval is covered by the selected original
bands; merging cannot enlarge the adopted union. For each parent row
a finite path of child vertices is checked. Each child interval is
nonempty; consecutive intervals intersect, either through both
cross-endpoint inequalities or through a checked common parent knot.
The first interval reaches the parent's lower endpoint and the last
reaches its upper endpoint. The union of a chain of intersecting closed
intervals is an interval, and these endpoint conditions imply that it
contains the parent interval. A point in a merged child band is then
assigned to a selected original band and a verified destination bin.
All extensions make positive total progress. Semantic checks validate
states, parity, endpoint identities, shape inclusions, and the ratio
bounds, so that the finite path represents actual extended prefixes.
Thus acceptance proves the successor property in Lemma~\ref{lem:cover}.

The complete replay checks every original adopted row, without deleting
rows that fail. Its recorded counts are
\[
\begin{array}{lr}
\text{adopted rows}&3,464,816\\
\text{adopted cells}&150,040\\
\text{selected child vertices}&799,932\\
\text{coverage-path vertices (with repetitions)}&6,541,897.
\end{array}
\]
The generators and exporter merely propose witnesses. Their output
must pass the semantic and coverage checks. The input algebraic endpoint
enclosures use denominator $2^{48}$; the original verifier regenerates
the input and checks byte-for-byte agreement with the saved file.
The Lean checker uses exact rational arithmetic and separately validates
the exact $\Q(\sqrt{462})$ labels underlying the enclosures.
The family has uniform positive lower and upper ratio bounds
(the formal real states use $1/10000\le S\le10000$), so
Lemma~\ref{lem:cover} applies.

\section{Attaching the root}
Take $A=322$ and $B=431$. Direct matrix multiplication gives
\[
 F_A(x)=\frac{5+2x}{17+7x},\qquad F_B(x)=\frac{4+3x}{17+13x}.
\]
Both words have odd length; their states are $0$ and $2$.
Lemma~\ref{lem:extrema} and substitution in these maps give
\[
 L=L_A+L_B=\frac{8512493+28004\dd}{17339617},\qquad
 H=H_A+H_B=\frac{1352829-8081\dd}{2233974}.
\]
The root shapes and scale ratio satisfy
\[
 r_A=7/17\in[9/22,14/33],\quad
 r_B=13/17\in[13/17,19/24],
\]
\[
 S(A,B)=\frac{55792801+2467176\dd}{159491641}
 \in[\kappa^{-20},\kappa^{-19}].
\]
The adopted root rows are the 25 bands
$(p,q)=(j/32,(j+2)/32)$, $j=2,\ldots,26$.
They are types numbered 3--27 at state pair 193 and ratio bin 135
in the stored format. Their union is exactly $J_{1/16,7/8}(A,B)$:
the first left coordinate is $2/32$, the last right coordinate is
$28/32$, and successive bands overlap. The header's other root index
is not used in place of these rows.
Applying the closed-cover principle to each of them gives
\begin{equation}\label{eq:sum}
 J_{1/16,7/8}(322,431)\subset K(322)+K(431).
\end{equation}
Set
\[
 \alpha=4+L+(H-L)/16,\qquad \beta=4+L+7(H-L)/8.
\]
They are approximately $4.525777278415714\ldots$ and
$4.527546990114258\ldots$. All comparisons used below are exact:
\begin{equation}\label{eq:ordering}
 \frac{4525423}{1000000}<\alpha<\frac{226289}{50000}
 <\frac{226377}{50000}<\beta<c_F.
\end{equation}
For an explicit rational method to check these inequalities, use
\[
 \frac{214941852602}{10^{10}}<\sqrt{462}
 <\frac{214941852603}{10^{10}}.
\]
Squaring the positive endpoints proves these bounds. Substitution into
each affine expression in $\sqrt{462}$, choosing the endpoint according
to the coefficient's sign, proves \eqref{eq:ordering} by integer
arithmetic. The same bounds check $L<H$ and the root ratio box.

\section{Uniform control at all noncentral positions}\label{sec:spectral}
The sum realization alone does not prove membership in $M$.
For $x\in K(322)$ and $y\in K(431)$ glue the corresponding digit
streams around a central $4$. The fixed core is
\[
 \cdots\,2\,2\,3\,4^*\,4\,3\,1\,\cdots,
\]
where the star marks position zero; the streams are read outwards.
We give the rational bound corresponding directly to the Lean proof.

For a finite word $w$ and a rational interval $I$, let $E(w,I)=F_w(I)$
with endpoints reordered, and write $E^+(w,I)$ for its upper endpoint.
This is computed recursively by
\[
 E(\varnothing,I)=I,\qquad
 E(a w,I)=\left[\frac1{a+E^+(w,I)},\frac1{a+E^-(w,I)}\right].
\]
All terminal intervals below are nonnegative, so no denominator vanishes.
Define rational outer intervals $R_s$ by
\[
\begin{array}{c|cc}
 s&\inf R_s&\sup R_s\\\hline
0&263931509/10^9&197795529/250000000\\
1&263931509/10^9&788861617/10^9\\
2&267649463/10^9&197795529/250000000\\
3&263931509/10^9&368115019/500000000\\
4&358271131/10^9&197795529/250000000
\end{array}
\]
The preceding rational square-root bounds show
$[l_s,u_s]\subset R_s$. Hence actual outward tails in state $s$ lie in $R_s$.
Put $b=4525423/1000000$.

\begin{lemma}[Finite rational spectral checks]\label{lem:finite}
The following finite inequalities hold.
\begin{enumerate}
\item For every $w\in\{1,2,3\}^6$ and $a\in\{1,2,3\}$ such that
$wa$ is legal from state zero, ending in state $s$,
\[
 a+E^+(w^{\mathrm{rev}},[0,1])+\sup R_s\le b.
\]
\item For $(P,Q,s_Q)=(322,431,2)$ or $(431,322,0)$,
for $0\le n<6$, $w\in\{1,2,3\}^n$, and $a\in\{1,2,3\}$,
if $Pwa$ is legal and ends in state $s$,
\[
 a+E^+((Pw)^{\mathrm{rev}}\,4\,Q,R_{s_Q})+\sup R_s\le b.
\]
\item At the adjacent fixed $4$,
\[
 4+E^+(4\,322,R_0)+E^+(31,R_2)\le b.
\]
\end{enumerate}
These finite assertions are checked by rational reduction inside the
Lean kernel in \file{SpectralVerified.lean}.
\end{lemma}
\begin{proof}
The enumerations and the displayed recurrence completely specify the
checks. There are $3^6\cdot3=2187$ cases in item 1,
$2\cdot3\sum_{n=0}^5 3^n=2184$ in item 2, and one in item 3,
for a total of 4372. Illegal cases are discarded by automaton traversal.
For each retained case recursively compute the two endpoints of $E$,
add the displayed rational upper bounds, and compare with $b$ by
cross multiplication of positive denominators. The four Boolean
acceptance lemmas are \file{far_checked}, \file{near322_checked},
\file{near431_checked}, and \file{adjacent_four_checked}; each is proved
with \file{decide +kernel}. This finite computation is the proof of
these rational inequalities, rather than a numerical estimate or an
unverified assertion about their maximum.
\end{proof}

\begin{proposition}\label{prop:noncentral}
For every glued pair of legal infinite streams,
$\lambda_n(\theta)\le b$ for every integer $n\ne0$.
\end{proposition}
\begin{proof}
Index either outward stream by $u_0,u_1,\ldots$, so the fixed prefix
occupies indices $0,1,2$, and the digit $u_k$ is at distance $k+1$
from the central position. For $3\le k<9$, write the stream up to
$u_k$ as $Pwa$, where $|w|=k-3<6$. Its inward continued fraction
starts with $(Pw)^{\mathrm{rev}}4Q$ and then a tail in $R_{s_Q}$.
Its outward tail lies in $R_s$, with $s$ the state after $Pwa$.
Item 2 of Lemma~\ref{lem:finite} bounds their sum with the digit $a$.

For $k\ge9$, retain the six digits $u_{k-6},\ldots,u_{k-1}=w$
and put $a=u_k$. They all lie in $\{1,2,3\}$. The inward fraction
begins with $w^{\mathrm{rev}}$; its remaining value lies in $[0,1]$,
even though it eventually crosses the fixed core. The outward fraction
lies in $R_s$. The actual state after $wa$ equals the state of the
original stream at $k+1$: the forbidden-word automaton depends only on
the last four digits. Thus item 1 bounds the actual local value.
This treats infinitely many positions with one finite list of checks,
without replacing the original tail by another sequence.

In the fixed core, the three noncentral digits at positions $-3,-2,3$ have local value at most $4$, because
both fractional tails lie in $[0,1]$. Positions $-1$ and $2$ have digit
$3$ and an adjacent digit at least $2$; the corresponding fractional
tail is at most $1/2$, the other at most $1$, so their values are at
most $9/2$. The remaining position $1$, with digit $4$, is covered
by item 3. Since $4\le9/2<b$, all positions are covered.
\end{proof}

The original exact-algebraic verifier also proves the sharper bound
$4\sqrt{462}/19+1/3025=4.525422212239698\ldots$, using legal five-digit
windows, nearby-prefix enumeration, and a nine-digit cylinder replacement.
That argument is additional evidence; the present proof needs only
Proposition~\ref{prop:noncentral}. The separate rational implementation
uses 40-digit extremal prefixes with terminal $[0,1]$; its enumeration
includes 242 legal five-digit windows and 57,054 nearby prefixes.
These counts differ from the 4372 checks above because the algorithms
differ, not because they verify different families of infinite tails.

\begin{proof}[Proof of Theorem~\ref{thm:main}]
Fix $t\in[\alpha,\beta]$. By \eqref{eq:sum}, choose $x\in K(322)$,
$y\in K(431)$ with $t=4+x+y$, and glue their digit streams.
Then $\lambda_0(\theta)=t$, whereas
$\lambda_n(\theta)\le b<\alpha\le t$ for $n\ne0$ by
Proposition~\ref{prop:noncentral} and \eqref{eq:ordering}.
Therefore $m(\theta)=t$ and $[\alpha,\beta]\subset M$.
Equation~\eqref{eq:ordering} gives the asserted rational subinterval
and its separation from $c_F$.
Every point of its open interior has a sufficiently small open
neighborhood contained in that interval and below $c_F$, proving the
interior statement. Subtraction and averaging of the exact rational
endpoints give the width and midpoint.
\end{proof}

\section{Formal verification and execution trust boundary}
The public theorem \file{Berstein.target_interval_subset} in
\file{Berstein/Lean/Berstein/Main.lean} has the following mathematical form:
\begin{gather*}
 \operatorname{Certificate.check}(\mathrm{input},\mathrm{data},
 \mathrm{alive},\mathrm{cells})=\mathrm{true}
 \\ \Longrightarrow \\
 [\mathrm{targetLower},\mathrm{targetUpper}]
 \subset \mathrm{markovSpectrum}\cap(-\infty,\mathrm{freimanConstant}).
\end{gather*}
Its arguments are \file{GraphCertificate.Input}, \file{GraphMeaning.Data},
a \file{ByteArray}, and a function assigning each geometry its list of
\file{GraphCertificate.Cell} witnesses. Its sole premise is finite-checker
acceptance; there are no separate assumed Cantor-sum, infinite-realization,
or noncentral-bound hypotheses.
The Boolean predicate is the conjunction of semantic input checking,
root checking, preparation checking on each side, and geometry checking
for all state pairs. The module correspondence is
\begin{center}\small
\begin{tabular}{p{.33\linewidth}p{.59\linewidth}}\toprule
Argument & Modules under \file{Berstein/Lean/Berstein/}\\\midrule
Infinite values and language & \file{Markov}, \file{Forbidden31313}, \file{CFInvariant}\\
Exact algebraic meanings & \file{Quadratic462}, \file{QuadraticCF}, \file{SemanticCheck}, \file{SemanticSound}\\
Box comparisons & \file{Normalization}, \file{SemanticTransport}, \file{GraphNumericSound}\\
Finite paths and closure & \file{GraphCertificate}, \file{GraphCertificateFacts}, \file{GraphSuccessor}\\
Infinite realization & \file{ActualCylinders}, \file{CylinderCover}, \file{GraphState}\\
Root and dominance & \file{RootFacts}, \file{RootCheck}, \file{RootState}, \file{SpectralVerified}\\
Final implication & \file{CertificateCheck}, \file{Verified}, \file{Main}\\\bottomrule
\end{tabular}
\end{center}
The formalization retains actual prefixes in addition to finite indices;
a verified graph cycle is not itself mistaken for an infinite
continued-fraction realization.

The Lean kernel checks this implication and the small spectral
computations. The concrete 3,464,816-row acceptance is established by
compiled execution of \file{graphReplay}, with four shards covering
all 484 geometries. The executable validates input semantics and roots,
prepared comparisons and every geometry. It is not a kernel reduction
producing a closed theorem for this concrete byte array. No external
success flag is postulated as an axiom. Applying the soundness theorem
to this concrete execution therefore trusts the compiler, runtime,
arbitrary-precision arithmetic implementation, file parsing and input
identity, and execution environment in addition to the kernel and
standard foundational axioms.

\file{Audit.lean} prints the axioms of the final theorems. The verification
script accepts only \file{propext}, \file{Classical.choice}, and
\file{Quot.sound}, the standard Lean/Mathlib axioms. The development
uses no \file{sorry}, \file{admit}, or \file{native_decide} to establish
these results. Hashes identify data, executable, witnesses, and logs;
they are evidence of identity, not substitutes for the soundness proof.
Exporters and search heuristics are untrusted: malformed proposals must
be rejected by the proved checker. Floating-point candidate selection
is not used for final acceptance.

The independent audit is a separate implementation in the same
repository, not external peer review. Its input checker reconstructs
quadratic values and transitions without importing the original Python
modules. Its C++ checker uses its own interval and connectivity
implementation to check every adopted row, with zero failed cells and
no row deletion in the recorded run. Its spectral checker uses rational
outer enclosures. Negative controls exercise gaps, missing endpoints,
missing child states or ratio boxes, and truncated data. There are also
4000 arithmetic comparisons against Python's \file{Fraction}.
The Lean corruption tests cover a root adoption bit, exact endpoint,
used constant, cell, coverage path, and destination mask. These checks
help detect execution or implementation errors; they do not remove the
compiler/runtime boundary.

\section{Reproduction and provenance}
The immutable release dated October 8, 2026 is cited in \cite{release}.
The paper is a subsequent exposition of that construction, not a
replacement of its immutable artifacts. To inspect the released source,
check out the release tag; to inspect the source used for the module
mapping in this paper, check out the snapshot commit stated above.
Run commands from the repository root unless specified otherwise.
A Python 3 interpreter and a C++17 compiler suffice for the original
and independent checkers. Do not run the audit with Python's \file{-O}
option, which disables assertions.
\begin{verbatim}
git clone https://github.com/K-Yamada-2002/verify-freiman-constant.git
cd verify-freiman-constant
git checkout markov-interior-2026-10-08
python3 -B -S Berstein/src/verify_below_ray.py --full
\end{verbatim}
This regenerates the algebraic input, requires byte-for-byte agreement,
and rechecks every adopted row without changing the original input or
adoption bitmap. Outputs include
\file{Berstein/data/below_ray_verified.json},
\file{below_ray_certificate.json}, and
\file{below_ray_kernel_replay.json}. Omitting \file{--full} reuses recorded
verification after hash checks; it is not a fresh exhaustive replay.
The independent audit commands are
\begin{verbatim}
python3 -I -B -S Berstein/audit/independent_input.py
python3 -I -B -S Berstein/audit/independent_spectral.py
c++ -std=c++17 -O2 -fsanitize=undefined -fno-sanitize-recover=all \
  Berstein/audit/independent_kernel.cpp -o /tmp/berstein-independent-kernel
/tmp/berstein-independent-kernel Freiman/data/graph_wide.dat \
  Freiman/data/graph_wide.json.alive.bin
python3 -I -B -S Berstein/audit/check_independent_kernel.py \
  /tmp/berstein-independent-kernel
python3 -I -B -S Berstein/audit/check_independent_arithmetic.py
\end{verbatim}
For the formal proof and exhaustive compiled replay, Lean 4.32.1 and
Mathlib v4.32.1 are pinned:
\begin{verbatim}
cd Berstein/Lean
lake update
lake exe cache get
python3 scripts/verify.py --prepare
\end{verbatim}
The final command generates witnesses, builds the proof and executable,
audits axioms, replays the entire table, and runs corruption tests.
With generated witnesses, \file{python3 scripts/verify.py} performs
verification without regeneration. The \file{--reuse-replay} option
reuses a successful replay only after checking hashes of all 484
certificates, original data, executable and logs; it reruns the build,
axiom audit and corruption tests. \file{lake build} alone does not replay
the table. \file{lake env lean Audit.lean} prints the axiom audit.
The committed historical record is
\file{Berstein/Lean/verification-baseline.json}; fresh combined results
are written to \file{Berstein/Lean/logs/verification.json}.
The recorded shard row totals are 869449, 868886, 871412 and 855069,
whose sum is 3464816. All recorded shard exit codes are zero.
The baseline records evidence from prior runs; it is not a claim that
preparing this manuscript reran the entire large computation.

Some continued-fraction modules were vendored from the author's
\file{hall-ray} repository at commit
\file{a584e5620124d23c709e3bbdc66db8b8b0586310}.
Their source identities and license are recorded in
\file{Berstein/Lean/sources.json} and \file{HallRay/LICENSE}.
No adjacent checkout is required at runtime.
The original Japanese exposition is
\file{Berstein/doc/generalized_t.tex}; the present article isolates its
completed Markov-interval construction and supplies a proof matching
the formal spectral bound.

\section{What is and is not concluded}
The completed claim is the displayed interval inclusion in $M$ and
its consequence for interior in the real line below $c_F$.
This proof does not establish inclusion in the Lagrange spectrum,
nor in $M\setminus L$. The $131$-forbidden problem discussed elsewhere
in the repository is different: $31313$-avoidance allows occurrences
of $131$. Its unfinished searches and certificates under
\file{stashed/} are not premises of this proof. Neither a local
one-step cover nor unsuccessful finite searches settle that separate
interior problem. Likewise, the classical identification of the
continued-fraction spectrum with the quadratic-form spectrum and
Freiman's maximal-ray theorem are cited background results, distinct
from the formally checked theorem about the explicitly defined
continued-fraction spectrum and constant.

\begin{thebibliography}{9}
\bibitem{CF} T. W. Cusick and M. E. Flahive,
\emph{The Markoff and Lagrange Spectra}, Mathematical Surveys and
Monographs 30, American Mathematical Society, 1989.
\url{https://doi.org/10.1090/surv/030}.
\bibitem{Freiman} G. A. Freiman,
\emph{Diophantine Approximations and the Geometry of Numbers
(Markov's Problem)}, Kalinin State University, Kalinin, 1975
(in Russian).
\bibitem{release} K. Yamada,
\emph{Interior points of the Markov spectrum outside Hall's ray},
immutable source release, October 8, 2026.
\url{https://github.com/K-Yamada-2002/verify-freiman-constant/releases/tag/markov-interior-2026-10-08}.
\end{thebibliography}
\end{document}
